\documentclass[11pt,letterpaper]{article}
\usepackage{bes-whitepaper}

\begin{document}

\wptitle{One barrier or two steps? Nucleation statistics for gypsum and\\how additives change its pathway}
  {Synthesis and Processing Science}{Shawn Chen}

\hd{Research question.} When gypsum (CaSO$_4\cdot$2H$_2$O) crystallizes from water, is the rate set by
a single nucleation barrier, as classical nucleation theory assumes, or by forming and then
transforming a precursor? And when an additive delays crystallization, which step does it act on?

\hd{Knowledge gap.} Gypsum is a model for low-energy crystallization of an ionic solid from water,
and a costly scale in geothermal, oil and gas, and water systems. Scattering and imaging show
nanoscale precursors [1,2] but cannot say whether the precursor step limits the rate. Nucleation is
stochastic, so that answer lies in the distribution of induction times over many identical samples
[3,4], yet scale and additive studies usually report one onset time per condition.

\hd{Hypotheses.} \emph{H1, single barrier:} each well crystallizes by one rare event, so the
fraction of wells still clear decays exponentially after a fixed growth lag, $\ln J$ falls linearly
with $1/\ln^2 S$, and an additive moves only the slope (interfacial energy) or intercept (prefactor) of that line.
\emph{H2, precursor-mediated:} two slow steps in series give a sigmoidal, non-exponential
distribution, a break in slope where the slower step switches, and additives that change the
distribution's shape, not just its scale.

\hd{Approach.} An Opentrons OT-2 fills 96-well plates from filtered CaCl$_2$ and Na$_2$SO$_4$
stocks to set the supersaturation $S$ (from Pitzer activities), ionic strength and additive (a
phosphonate, a polyacrylate, citrate), with 48--96 replicate wells per condition and a logged start
time for each well. Turbidity is read every 30 s by a plate reader, or by the lab's camera under
fixed lighting. Fill volume is varied too: rates that scale with volume mean nucleation in the bulk,
and rates that scale with wetted area mean the walls. Selected wells are sampled at set times to
identify the first solid (Raman, XRD). Replicates are added where H1 and H2 predict the most
different distributions; the automation is the means, not the aim.

\wpflow
  {\textbf{Design}\\ $S$, additive and fill volume; add replicates where H1, H2 differ}
  {\textbf{OT-2}\\ 48--96 replicate wells per condition; per-well start times}
  {\textbf{Turbidity}\\ every 30 s; first solid by Raman or XRD}
  {\textbf{Statistics}\\ induction-time distributions; fit H1 and H2}
  {\textbf{Outcome}\\ rate-limiting step; where each additive acts}
  {update the probability of H1 and H2, then choose the next plate}

\hd{Expected outcomes.} (1) Nucleation rates with uncertainties across $S$, ionic strength and
additives, with bulk and wall contributions separated; (2) a test of whether precursor formation
controls the rate; (3) which step each additive acts on. Together these give a statistically
grounded time--supersaturation--transformation map for a model ionic crystal, and a mechanistic basis
for choosing additives.

\hd{Fit and resources.} The project asks why crystallization proceeds as it does, through the roles
of \emph{nucleation} and \emph{precursor and intermediate states} in materials synthesis. It does not
optimize a process or develop automation for its own sake. The PI works on self-driving labs and
Bayesian experimental design; the lab's OT-2 already runs turbidity-onset precipitation series, so
preliminary distributions can be measured before an application.

\begin{wprefs}
[1] A.\,E.\,S. Van Driessche \emph{et al.}, \emph{Science} 336, 69 (2012), doi:10.1126/science.1215648.
[2] T.\,M. Stawski \emph{et al.}, \emph{Nat.\ Commun.} 7, 11177 (2016), doi:10.1038/ncomms11177.
[3] S. Jiang, J.\,H. ter Horst, \emph{Cryst.\ Growth Des.} 11, 256 (2011), doi:10.1021/cg101213q.
[4] T. Brandel, J.\,H. ter Horst, \emph{Faraday Discuss.} 179, 199 (2015), doi:10.1039/C4FD00230J.
\end{wprefs}

\end{document}
